% !TeX root = ../main.tex
\section{Results}
\label{sec:results}

\subsection{RQ1: Applicability-Classification Outcomes}
\label{sec:rq1-results}

The exhaustive assessment produced \StudyAssessmentCount{} unique
story--pattern outcomes. Of these, \StudyApplicableCount{}
(\RqOneApplicablePercent{}\%) were classified as applicable,
\StudyNotApplicableCount{} (\RqOneNotApplicablePercent{}\%) as not applicable,
and \StudyInsufficientInformationCount{}
(\RqOneInsufficientInformationPercent{}\%) as insufficient information
(Table~\ref{tab:applicability-results}). Each percentage uses all assessed
pairs as its denominator. These are model-label frequencies, not estimates
of classification accuracy or expert-confirmed applicability.

% !TeX root = ../main.tex
% Generated by processing/scripts/prepare_applicability_results.py.
\begin{table}[t]
\caption{RQ1 model-classification outcomes. Percentages use assessed pairs in the indicated row, not expert judgments.}
\label{tab:applicability-results}
\centering
\small
\begin{tabular}{lrrr}
\hline
\multicolumn{4}{l}{\textbf{Overall classification distribution}} \\
Classification & Count & Pairs & \% \\
\hline
Applicable & \StudyApplicableCount{} & \StudyAssessmentCount{} & \RqOneApplicablePercent{} \\
Not applicable & \StudyNotApplicableCount{} & \StudyAssessmentCount{} & \RqOneNotApplicablePercent{} \\
Insufficient information & \StudyInsufficientInformationCount{} & \StudyAssessmentCount{} & \RqOneInsufficientInformationPercent{} \\
\hline
\multicolumn{4}{l}{\textbf{Applicable outcomes by story and system}} \\
Story / subtotal & Applicable & Pairs & \% \\
\hline
SYS01-US01 & 3 & \StudyPatternCount{} & 4.29 \\
SYS01-US02 & 1 & \StudyPatternCount{} & 1.43 \\
SYS01-US03 & 3 & \StudyPatternCount{} & 4.29 \\
\textbf{SYS01 subtotal} & \RqOneSystemOneApplicableCount{} & 210 & 3.33 \\
\hline
SYS02-US01 & 5 & \StudyPatternCount{} & 7.14 \\
SYS02-US02 & 5 & \StudyPatternCount{} & 7.14 \\
SYS02-US03 & 2 & \StudyPatternCount{} & 2.86 \\
\textbf{SYS02 subtotal} & \RqOneSystemTwoApplicableCount{} & 210 & 5.71 \\
\hline
SYS03-US01 & 1 & \StudyPatternCount{} & 1.43 \\
SYS03-US02 & 5 & \StudyPatternCount{} & 7.14 \\
SYS03-US03 & 2 & \StudyPatternCount{} & 2.86 \\
\textbf{SYS03 subtotal} & \RqOneSystemThreeApplicableCount{} & 210 & 3.81 \\
\hline
SYS04-US01 & 0 & \StudyPatternCount{} & 0.00 \\
SYS04-US02 & 2 & \StudyPatternCount{} & 2.86 \\
SYS04-US03 & 7 & \StudyPatternCount{} & 10.00 \\
\textbf{SYS04 subtotal} & \RqOneSystemFourApplicableCount{} & 210 & 4.29 \\
\hline
All systems & \StudyApplicableCount{} & \StudyAssessmentCount{} & \RqOneApplicablePercent{} \\
\hline
\end{tabular}
\end{table}


Applicable outcomes occurred in \StudyPositiveStoryCount{} of
\StudyStoryCount{} stories. Counts ranged from \RqOneStoryMinimum{} to
\RqOneStoryMaximum{} per story, with each story assessed against
\StudyPatternCount{} patterns. \texttt{SYS04-US01} had no applicable
outcome; this is a retained classification result, not an execution failure
or evidence that no pattern could apply in practice.
The largest positive set belonged to \texttt{\RqOneMaximumStories{}},
with \RqOneStoryMaximum{} applicable outcomes.

At system level, applicable counts were \RqOneSystemOneApplicableCount{},
\RqOneSystemTwoApplicableCount{}, \RqOneSystemThreeApplicableCount{}, and
\RqOneSystemFourApplicableCount{} for SYS01--SYS04, respectively.
System percentages in the table use the assessed pairs within that system,
rather than the total positive set. The insufficient-information category
accounted for most outcomes overall. The observed differences describe this
bounded sample; they do not identify causes, establish that the requirements
are defective, or imply that non-positive outcomes are missed applicable
patterns. Architectural adequacy requires separate evidence.

\subsection{RQ2: Overlap with the Retrieval Baseline}
\label{sec:rq2-results}

Of the \StudyApplicableCount{} LLM-positive story--pattern pairs,
\StudyRagOverlapCount{} appeared in the frozen baseline's final Top-5 for
the same story, yielding a pooled overlap of \RqTwoOverlapPercent{}\%
(Equation~\ref{eq:rag-overlap}). These matches spanned
\StudyRagOverlapStoryCount{} of \StudyStoryCount{} sampled stories.
The remaining \StudyRagOutsideCount{} positive pairs were outside that
story's final Top-5; they are not established false negatives.

% !TeX root = ../main.tex
% Generated by processing/scripts/prepare_rag_results.py.
\begin{table}[t]
\caption{RQ2 final Top-5 membership among LLM-positive pairs. Percentage denominator: positive pairs in the row; -- denotes an undefined percentage for a zero-positive story.}
\label{tab:rag-comparison}
\centering
\small
\begin{tabular}{lrrrr}
\hline
Story & Positive & In & Out & \% in \\
\hline
SYS01-US01 & 3 & 1 & 2 & 33.33 \\
SYS01-US02 & 1 & 1 & 0 & 100.00 \\
SYS01-US03 & 3 & 0 & 3 & 0.00 \\
SYS02-US01 & 5 & 1 & 4 & 20.00 \\
SYS02-US02 & 5 & 1 & 4 & 20.00 \\
SYS02-US03 & 2 & 1 & 1 & 50.00 \\
SYS03-US01 & 1 & 0 & 1 & 0.00 \\
SYS03-US02 & 5 & 0 & 5 & 0.00 \\
SYS03-US03 & 2 & 0 & 2 & 0.00 \\
SYS04-US01 & 0 & 0 & 0 & -- \\
SYS04-US02 & 2 & 0 & 2 & 0.00 \\
SYS04-US03 & 7 & 0 & 7 & 0.00 \\
\hline
Pooled & \StudyApplicableCount{} & \StudyRagOverlapCount{} & \StudyRagOutsideCount{} & \RqTwoOverlapPercent{} \\
\hline
\end{tabular}
\end{table}


Table~\ref{tab:rag-comparison} retains every sampled story and reports
the positive-set size, pairs inside and outside the final Top-5, and
the overlap percentage among that story's LLM-positive pairs.
For \texttt{SYS04-US01}, the positive set is empty: the intersection count
is zero, but its percentage is undefined, not zero.
The pooled percentage divides total matches by all positive pairs;
it is not the average of the story percentages or a proportion of
retrieval candidates.

This comparison concerns final Top-5 membership only; vector Top-15
membership, ranking positions, and scores do not contribute to this measure.
The LLM labels are predictions rather than an independent reference set,
so the intersection does not measure precision, recall, classification
accuracy, or superiority of either approach. As documented in
Section~\ref{sec:rag-baseline}, the baseline uses English queries and a
richer catalog representation than the Portuguese-input applicability
assessment. Consequently, observed overlap differences cannot be
attributed solely to the selection mechanism.

\subsection{ADR-Generation Completion and Provenance}
\label{sec:adr-results}

The frozen generation outputs contain \StudyAdrCount{} unique ADRs for the
\StudyApplicableCount{} LLM-positive story--pattern pairs. The output pair
set equals the selected-pair set: every selected pair has exactly one final
ADR, with no missing or additional pair. The approved
\texttt{SYS01-US01}--Pattern~\texttt{001} pilot is included in this total,
not counted again as a separate observation.

All \StudyAdrCount{} records contain nonempty string-valued title, context,
decision, and consequences fields and the fixed status \texttt{Proposed}.
The offline validation checks exact field membership, pair identifiers,
source and rendered-prompt hashes, model configuration, response metadata,
and stored strict-output validation flags. Recompiling the frozen records
reproduces both the ADR-field CSV and the ADR-enriched selected-pair CSV
byte-for-byte. The source hash refers to the preserved pre-ADR selection
file, rather than its subsequently enriched version.

The execution history retains a \emph{Connection error} for
\texttt{SYS04-US03}--Pattern~\texttt{003}, followed by an explicitly
authorized retry linked to the original attempt. The failed attempt and
retry authorization are execution metadata, not additional ADRs;
the final record for that pair is counted once. The connection failure
does not establish whether the failed request reached the provider, so
the completed-record count is not a verified total of all API dispatches.

These results establish artifact availability, field completeness, and
traceable one-to-one generation for the selected pairs. Nonempty fields,
schema conformity, and reproducible compilation do not demonstrate that
the decisions are correct, sufficiently grounded, useful, or accepted by
architects. \texttt{Proposed} is an application-assigned status, not an
expert endorsement. Quality judgments are reserved for the separate
expert assessment; no such finding is asserted here.
